\begin{textsample}{POS Dim 3 – ro – Score 23.00 – ro/ro\_inf\_26.txt}  \label{ex:f3_pos_124}
Corpo, \textbf{gesto} e movimentos \textbf{Crianças} \textbf{pequenas} ( 4 anos a 5 anos e 11 \textbf{meses} ) - Interagem com outras \textbf{crianças} em \textbf{brincadeiras} de faz-de-conta, atividades de culinária, de manipulação de argila ou de manutenção de uma horta, de reconto coletivo de história, de construção com sucata, de \textbf{pintura} coletiva de um cartaz. - Participem de jogos de regras e aprendem a construir estratégias de jogo. - Arrumem a mesa para um almoço com os amigos, e manter a organização de seus pertences. - Ouvem e recontam histórias dos povos indígenas, africanos, asiáticos, europeus, de diferentes regiões do Brasil e de outros países da América. - Localizem em um mapa, com apoio do/a professor/a, sua cidade, ou aldeia, ou assentamento, e o local do Brasil no mapa mundial. - Participem de rodas de conversa para falar de situações pessoais ou narrar histórias familiares no grupo e ser ouvido por todos. - Discutem em grupo situações-problemas ou formas de planejar um evento. - Preparem uma exposição de objetos relativos às atividades e profissões dos familiares e dos \textbf{adultos} da unidade de Educação \textbf{Infantil}. - Pesquisem em casa suas tradições familiares, de modo a reconhecer elementos da sua identidade cultural. - Estabelecem relações entre o modo de vida característico de seu grupo social e o de outros grupos.
Conhecem costumes e \textbf{brincadeiras} de outras épocas e de outras civilizações. - Explorem \textbf{brincadeiras}, tipos de alimentação e de organização social característicos de diferentes culturas. - Explorem fotografias do seu grupo de \textbf{crianças} em diversas situações. - Experimentem sabores, percebem cheiros dos alimentos e escolhem o que querem comer. - Observem o ambiente e percebem aromas, \textbf{texturas}, sonoridades na companhia de outras \textbf{crianças}. - Comentem em diferentes linguagens com a \textbf{professora} sobre suas fotos e as de seus familiares. - Ouvem histórias lidas ou \textbf{contadas} pela \textbf{professora} e cantam com ela e as demais \textbf{crianças}. - \textbf{Brinquem} diante do espelho, observando os próprios \textbf{gestos} ou \textbf{imitando} outras \textbf{crianças}. - Participem junto com outras \textbf{crianças} de \textbf{refeições} gostosas e cheirosas, de momentos de \textbf{descanso} \textbf{diário} em ambiente aconchegante e silencioso, de momentos de banho refrescante. - \textbf{Brinquem} de \textbf{esconder} -se, de cuidar de animais \textbf{domésticos}, de ouvir e \textbf{contar} histórias, observam aspectos do ambiente, colecionam objetos, participam de \textbf{brincadeiras} de roda, \textbf{brincam} de faz-de-conta, dentre outras experiências realizadas com diferentes parceiros. - Vestem fantasias, experimentando ser outras pessoas, ou personagens de histórias que lhes são \textbf{contadas} ou lidas. - Torcem a favor de um grupo : um time esportivo, uma equipe musical, um grupo de gincana. - Cantem, respeitando sua vez de cantar e ouvindo os companheiros. - Decidem com os companheiros o tema de uma história a ser por todos \textbf{dramatizada}, usando esclarecimentos, justificativas e argumentos, que são muito ligados a seus sentimentos. - \textbf{Apoiem} parceiros em dificuldade, sem discriminá -los por suas características. - Realizem com maior autonomia ações de escovar os dentes, calçar sapatos ou o agasalho, pentear os cabelos, servir -se sozinha nas \textbf{refeições}, utilizar talheres adequados, lavar as mãos antes das \textbf{refeições} e depois de usar \textbf{tinta} ou \textbf{brincar} com terra ou \textbf{areia}. - \textbf{Brinquem} no \textbf{pátio}, praça ou jardim, em constante contato com a natureza.
Fonte : MEC.
Campos de experiências : efetivando direitos e aprendizagens na educação \textbf{infantil} ; texto final Zilma de Moraes Ramos de Oliveira, 2018.

% matched lemmas: adulto, apoiar, areia, brincadeira, brincar, contar, criança, descanso, diário, doméstico, dramatizar, esconder, gesto, imitar, infantil, mês, pequeno, pintura, professora, pátio, refeição, textura, tinta
\end{textsample}
